% !TeX program = pdflatex
\documentclass[11pt,a4paper]{amsart}

\newcommand{\notelanguage}{finnish}
% Shared preamble for course notes and exercise sets.

\usepackage[T1]{fontenc}
\usepackage{lmodern}
\providecommand{\notelanguage}{english}
\usepackage[finnish,english]{babel}
\AtBeginDocument{\expandafter\selectlanguage\expandafter{\notelanguage}}

\newcommand{\notesfinnishlanguage}{finnish}
\ifx\notelanguage\notesfinnishlanguage
\newcommand{\notestheoremname}{Lause}
\newcommand{\noteslemmaname}{Lemma}
\newcommand{\notespropositionname}{Propositio}
\newcommand{\notescorollaryname}{Korollaari}
\newcommand{\notesdefinitionname}{Määritelmä}
\newcommand{\notesexamplename}{Esimerkki}
\newcommand{\notesproblemname}{Tehtävä}
\newcommand{\notesexercisename}{Tehtävä}
\newcommand{\notesremarkname}{Huomautus}
\newcommand{\notessolutionname}{Ratkaisu}
\else
\newcommand{\notestheoremname}{Theorem}
\newcommand{\noteslemmaname}{Lemma}
\newcommand{\notespropositionname}{Proposition}
\newcommand{\notescorollaryname}{Corollary}
\newcommand{\notesdefinitionname}{Definition}
\newcommand{\notesexamplename}{Example}
\newcommand{\notesproblemname}{Problem}
\newcommand{\notesexercisename}{Exercise}
\newcommand{\notesremarkname}{Remark}
\newcommand{\notessolutionname}{Solution}
\fi

\usepackage[a4paper,margin=30mm]{geometry}
\usepackage{microtype}
\usepackage{graphicx}
\usepackage{enumitem}

\newcommand{\exerciseheading}[2]{%
  \par\noindent
  {\large\bfseries #1\par}
  \vspace{0.1em}
  \noindent{\large #2\par}
  \vspace{1.5em}
}

\usepackage{amsmath}
\usepackage{amssymb}
\usepackage{amsthm}
\usepackage{mathtools}
\usepackage{aliascnt}

% Row operations beside the affected row, e.g. \rowop{+2R_1}.
\newcommand{\rowop}[1]{\mathrel{\big|}\,#1}

% Linear systems: the argument gives the number of alignment pairs.
\newenvironment{linearsystem}[1]
{\left\{
  \begin{alignedat}{#1}}
    {
  \end{alignedat}\right.}

  \usepackage[hidelinks,pdfusetitle]{hyperref}

  \numberwithin{equation}{section}
  \setlist{itemsep=0.25em,topsep=0.5em}

  \theoremstyle{plain}
  \newtheorem{theorem}{\notestheoremname}[section]
  \newaliascnt{lemma}{theorem}
  \newtheorem{lemma}[lemma]{\noteslemmaname}
  \aliascntresetthe{lemma}
  \newaliascnt{proposition}{theorem}
  \newtheorem{proposition}[proposition]{\notespropositionname}
  \aliascntresetthe{proposition}
  \newaliascnt{corollary}{theorem}
  \newtheorem{corollary}[corollary]{\notescorollaryname}
  \aliascntresetthe{corollary}

  \theoremstyle{definition}
  \newaliascnt{definition}{theorem}
  \newtheorem{definition}[definition]{\notesdefinitionname}
  \aliascntresetthe{definition}
  \newaliascnt{example}{theorem}
  \newtheorem{example}[example]{\notesexamplename}
  \aliascntresetthe{example}
  \newaliascnt{problem}{theorem}
  \newtheorem{problem}[problem]{\notesproblemname}
  \aliascntresetthe{problem}
  \newtheorem{exercise}{\notesexercisename}

  \theoremstyle{remark}
  \newaliascnt{remark}{theorem}
  \newtheorem{remark}[remark]{\notesremarkname}
  \aliascntresetthe{remark}

  \newenvironment{solution}
  {
  \begin{proof}[\notessolutionname]}
    {
  \end{proof}}

  \usepackage[nameinlink,noabbrev]{cleveref}

  \crefname{theorem}{theorem}{theorems}
  \crefname{lemma}{lemma}{lemmas}
  \crefname{proposition}{proposition}{propositions}
  \crefname{corollary}{corollary}{corollaries}
  \crefname{definition}{definition}{definitions}
  \crefname{example}{example}{examples}
  \crefname{problem}{problem}{problems}
  \crefname{exercise}{exercise}{exercises}
  \crefname{remark}{remark}{remarks}

  \ifx\notelanguage\notesfinnishlanguage
  \crefname{theorem}{lause}{lauseet}
  \crefname{lemma}{lemma}{lemmat}
  \crefname{proposition}{propositio}{propositiot}
  \crefname{corollary}{korollaari}{korollaarit}
  \crefname{definition}{määritelmä}{määritelmät}
  \crefname{example}{esimerkki}{esimerkit}
  \crefname{problem}{tehtävä}{tehtävät}
  \crefname{exercise}{tehtävä}{tehtävät}
  \crefname{remark}{huomautus}{huomautukset}
  \fi

  \newcommand{\NN}{\mathbb{N}}
  \newcommand{\NNone}{\mathbb{N}_1}
  \newcommand{\ZZ}{\mathbb{Z}}
  \newcommand{\QQ}{\mathbb{Q}}
  \newcommand{\RR}{\mathbb{R}}
  \newcommand{\CC}{\mathbb{C}}
  \newcommand{\PP}{\mathbb{P}}

  \DeclareMathOperator{\im}{im}
  \DeclareMathOperator{\rank}{rank}
  \DeclareMathOperator{\syt}{syt}
  \DeclareMathOperator{\pyj}{pyj}
  \DeclareMathOperator{\dom}{dom}
  \DeclarePairedDelimiter{\abs}{\lvert}{\rvert}
  \DeclarePairedDelimiter{\norm}{\lVert}{\rVert}
  \newcommand{\set}[1]{\left\lbrace#1\right\rbrace}
  \newcommand{\maxset}[1]{\max\set{#1}}
  \newcommand{\minset}[1]{\min\set{#1}}
  \newcommand{\powerset}[1]{\mathcal{P}\!\left(#1\right)}

\usepackage{tikz}
\usepackage{needspace}

\title{Analyysi II\\Harjoitus 3}
\author{}
\date{}

\begin{document}

\exerciseheading{Analyysi II}{Harjoitus 3}

\begin{exercise}
  Olkoon \(f \colon A \to \RR\) aidosti kasvava funktio. Osoita, että \(f\)
  on injektio.
\end{exercise}

\begin{solution}
  Olkoot \(x, y \in A\) mielivaltaisia siten, että \(x \neq y\).
  Tällöin joko \(x < y\) tai \(y < x\). Funktion \(f\) aidon kasvavuuden
  nojalla ensimmäisessä tapauksessa \(f(x) < f(y)\) ja toisessa
  \(f(y) < f(x)\). Molemmissa tapauksissa siis \(f(x) \neq f(y)\),
  joten \(f\) on injektio.
\end{solution}

\begin{exercise}
  Perustele kuvien avulla, että \(\cos(-x) = \cos x\),
  \(\sin(-x) = -\sin x\) ja \(\cos(\pi/2 - x) = \sin x\).
\end{exercise}

\begin{solution}
  Olkoon \(x \in \RR\). Yksikköympyrällä kulmaa \(x\) vastaava piste on
  \(P_x = (\cos x, \sin x)\), kun kulma mitataan positiivisesta
  vaaka-akselista vastapäivään.

  Peilaus vaaka-akselin suhteen muuttaa kulman \(x\) kulmaksi \(-x\),
  säilyttää vaakakoordinaatin ja vaihtaa pystykoordinaatin vastaluvukseen
  (vasen kuva). Siksi
  \[
    P_{-x} = (\cos x, -\sin x).
  \]
  Toisaalta \(P_{-x} = (\cos(-x), \sin(-x))\), joten koordinaatteja
  vertaamalla saadaan
  \[
    \cos(-x) = \cos x
    \qquad\text{ja}\qquad
    \sin(-x) = -\sin x.
  \]

  Peilaus suoran \(v = u\) suhteen vaihtaa pisteen koordinaatit keskenään
  ja muuttaa kulman \(x\) kulmaksi \(\pi/2 - x\) (oikea kuva). Näin ollen
  \[
    P_{\pi/2-x} = (\sin x, \cos x).
  \]
  Tämän pisteen vaakakoordinaatti on määritelmän mukaan
  \(\cos(\pi/2 - x)\), joten
  \[
    \cos(\pi/2 - x) = \sin x.
  \]

  \begin{center}
    \begin{minipage}[t]{0.48\linewidth}
      \centering
      \begin{tikzpicture}[scale=1.8,>=stealth,
        every node/.style={font=\small}]
        \draw[->] (-1.2,0) -- (1.3,0) node[right] {\(u\)};
        \draw[->] (0,-1.2) -- (0,1.25) node[above] {\(v\)};
        \draw (0,0) circle (1);
        \coordinate (P) at (35:1);
        \coordinate (Q) at (-35:1);
        \draw[thick] (P) -- (0,0) -- (Q);
        \draw[dashed] (0,0.574) -- (P) -- (Q) -- (0,-0.574);
        \draw[->] (0.4,0) arc[start angle=0,end angle=35,radius=0.4];
        \node at (18:0.58) {\(x\)};
        \draw[->] (0.4,0) arc[start angle=0,end angle=-35,radius=0.4];
        \node at (-18:0.61) {\(-x\)};
        \fill (P) circle (1pt) node[above right] {\(P_x\)};
        \fill (Q) circle (1pt) node[below right] {\(P_{-x}\)};
        \node[below left] at (0,0) {\(0\)};
      \end{tikzpicture}

      \small Peilaus vaaka-akselin suhteen.
    \end{minipage}\hfill
    \begin{minipage}[t]{0.48\linewidth}
      \centering
      \begin{tikzpicture}[scale=1.8,>=stealth,
        every node/.style={font=\small}]
        \draw[->] (-1.2,0) -- (1.3,0) node[right] {\(u\)};
        \draw[->] (0,-1.2) -- (0,1.25) node[above] {\(v\)};
        \draw (0,0) circle (1);
        \draw[dashed] (-0.85,-0.85) -- (1.05,1.05)
        node[above right] {\(v=u\)};
        \coordinate (P) at (25:1);
        \coordinate (Q) at (65:1);
        \draw[thick] (P) -- (0,0) -- (Q);
        \draw[dotted] (P) -- (Q);
        \draw[->] (0.36,0) arc[start angle=0,end angle=25,radius=0.36];
        \node at (12:0.53) {\(x\)};
        \draw[->] (0.7,0) arc[start angle=0,end angle=65,radius=0.7];
        \node[fill=white,inner sep=1pt] at (46:0.7)
        {\(\frac{\pi}{2}-x\)};
        \fill (P) circle (1pt) node[right] {\(P_x\)};
        \fill (Q) circle (1pt) node[above,xshift=8pt] {\(P_{\pi/2-x}\)};
        \node[below left] at (0,0) {\(0\)};
      \end{tikzpicture}

      \small Peilaus suoran \(v=u\) suhteen.
    \end{minipage}
  \end{center}

  Kuvissa \(0 < x < \pi/2\), mutta samat peilaukset toimivat kaikissa
  neljänneksissä ja myös koordinaattiakseleilla. Kulmat ymmärretään
  täysien kierrosten tarkkuudella, joten perustelut pätevät kaikilla
  \(x \in \RR\).
\end{solution}

\begin{exercise}
  Tarkastellaan funktiota \(f \colon (1, 4) \to \RR\),
  \[
    f(x) = \frac{1}{x - 1} + \frac{1}{x - 4}.
  \]
  \begin{enumerate}[label=(\alph*)]
    \item Osoita, että on olemassa sellainen piste \(z \in (1, 4)\), jolle
      \(f(z) = \frac{1}{\pi}\).
    \item Tutki, saavuttaako funktio \(f\) avoimella välillä \((1, 4)\)
      pienimmän tai suurimman arvonsa.
  \end{enumerate}
\end{exercise}

\Needspace{5\baselineskip}
\begin{solution}
  \leavevmode\par
  \begin{enumerate}[label=(\alph*), nosep]
    \item Vakiofunktio \(x \mapsto 1\) sekä funktiot \(x \mapsto x-1\)
      ja \(x \mapsto x-4\) ovat polynomifunktioina jatkuvia (lause 4.1.14).
      Koska \(x-1 \neq 0\) ja \(x-4 \neq 0\) kaikilla \(x \in (1, 4)\),
      osamäärän ja summan jatkuvuuden nojalla \(f\) on jatkuva välillä
      \((1, 4)\) (lause 4.1.12). Erityisesti \(f\) on jatkuva suljetulla
      välillä \([2, 3] \subset (1, 4)\).

      Huomataan, että
      \[
        f(2) = 1 - \frac{1}{2} = \frac{1}{2}
        \qquad\text{ja}\qquad
        f(3) = \frac{1}{2} - 1 = -\frac{1}{2}.
      \]
      Koska \(0 < 2 < \pi\), saadaan
      \[
        f(3) < \frac{1}{\pi} < f(2).
      \]
      Täten väliarvolauseen (Bolzanon lauseen korollaari) nojalla
      on olemassa \(z \in (2, 3) \subset (1, 4)\) siten, että
      \[
        f(z) = \frac{1}{\pi}.
      \]
    \item Olkoon \(M \in \RR\) mielivaltainen ja valitaan
      \(\delta = 1/(\abs{M}+1)\). Tällöin \(0 < \delta \leq 1\).
      Kun \(1 < x < 1+\delta \leq 2\), on \(1/(x-4) > -1/2\), joten
      \[
        f(x) > \frac{1}{x-1} - \frac12
        > \frac{1}{\delta} - \frac12
        = \abs{M} + \frac12 > M.
      \]
      Siis \(\lim_{x \to 1^+} f(x) = \infty\).

      Vastaavasti, kun \(3 \leq 4-\delta < x < 4\), on
      \(1/(x-1) < 1/2\), joten
      \[
        f(x) < \frac12 + \frac{1}{x-4}
        < \frac12 - \frac{1}{\delta}
        = -\abs{M} - \frac12 < M.
      \]
      Siis \(\lim_{x \to 4^-} f(x) = -\infty\).

      Oletetaan, että \(f\) saavuttaa suurimman arvonsa pisteessä
      \(t \in (1, 4)\). Ensimmäisen raja-arvon nojalla on kuitenkin olemassa
      \(x \in (1, 4)\), jolle \(f(x) > f(t)\). Tämä on ristiriita.

      Oletetaan sitten, että \(f\) saavuttaa pienimmän arvonsa pisteessä
      \(s \in (1, 4)\). Jälkimmäisen raja-arvon nojalla on kuitenkin olemassa
      \(x \in (1, 4)\), jolle \(f(x) < f(s)\).
      Tämäkin on ristiriita. Siis \(f\) ei saavuta suurinta eikä pienintä
      arvoaan välillä \((1, 4)\).
  \end{enumerate}
\end{solution}

\begin{exercise}
  Olkoon \(f \colon \RR \to \RR\),
  \[
    f(x) = \frac{\cos x}{1 + \abs{x}}.
  \]
  Osoita, että funktio \(f\) saavuttaa suurimman ja pienimmän arvonsa
  joukossa \(\RR\).
\end{exercise}

\begin{solution}
  Kosinifunktio on jatkuva joukossa \(\RR\). Epäyhtälö
  \[
    \bigl|\abs{x}-\abs{y}\bigr| \leq \abs{x-y}
    \qquad (x,y \in \RR)
  \]
  osoittaa, että myös itseisarvofunktio on jatkuva. Näin ollen
  \(x \mapsto 1+\abs{x}\) on jatkuva ja kaikkialla positiivinen,
  joten osamäärän jatkuvuuden nojalla \(f\) on jatkuva joukossa \(\RR\).

  Olkoon \(\varepsilon > 0\) ja valitaan
  \(H = \maxset{1,1/\varepsilon}\). Kun \(\abs{x} > H\), saadaan
  \[
    \abs{f(x)} = \frac{\abs{\cos x}}{1+\abs{x}}
    \leq \frac{1}{1+\abs{x}}
    < \frac{1}{\abs{x}} < \varepsilon.
  \]
  Täten \(\lim_{x \to \infty} f(x) = \lim_{x \to -\infty} f(x) = 0\).

  Huomataan, että
  \[
    f(0) = 1
    \qquad\text{ja}\qquad
    f(\pi) = -\frac{1}{1+\pi} < 0.
  \]
  Edellisten raja-arvojen nojalla voidaan valita \(m < 0\) ja \(M > \pi\)
  siten, että \(\abs{f(x)} < 1/(1+\pi)\), kun \(x < m\) tai \(x > M\).
  Näillä \(x\):n arvoilla siis
  \[
    f(\pi) = -\frac{1}{1+\pi}
    < f(x) < \frac{1}{1+\pi} < f(0).
  \]

  Koska \(f\) on jatkuva suljetulla ja rajoitetulla välillä \([m,M]\),
  Weierstrassin min-max-lauseen nojalla on olemassa \(a,b \in [m,M]\)
  siten, että
  \[
    f(a) \leq f(x) \leq f(b)
    \qquad\text{kaikilla } x \in [m,M].
  \]
  Lisäksi \(0,\pi \in [m,M]\), joten välin ulkopuolella pätee
  \[
    f(a) \leq f(\pi) < f(x) < f(0) \leq f(b).
  \]
  Siis \(f\) saavuttaa pienimmän arvonsa pisteessä \(a\) ja suurimman
  arvonsa pisteessä \(b\) koko joukossa \(\RR\).
\end{solution}

\begin{exercise}
  Olkoon \(f \colon (0, \infty) \to \RR\), \(f(x) = \frac{1}{x}\).
  Osoita, että funktio \(f\) on tasaisesti jatkuva jokaisella välillä
  \([r, \infty)\), \(r > 0\), mutta se ei ole tasaisesti jatkuva välillä
  \((0, \infty)\).
\end{exercise}

\begin{solution}
  Olkoon \(r > 0\) ja \(\varepsilon > 0\) mielivaltaisia.
  Valitaan \(\delta = r^2\varepsilon > 0\).
  Kun \(x,y \in [r, \infty)\), niin selvästi \(x,y \geq r\), joten
  \(xy \geq r^2\). Täten
  \[
    \abs{f(x) - f(y)} = \frac{\abs{x - y}}{xy} \leq \frac{\abs{x -
    y}}{r^2} < \varepsilon,
  \]
  kaikilla \(x,y \in [r, \infty)\), joilla \(\abs{x - y} < \delta\). Täten
  \(f\) on tasaisesti jatkuva välillä \([r, \infty)\).
  Koska \(r\) oli mielivaltainen, \(f\) on tasaisesti jatkuva välillä
  \([r, \infty)\) kaikilla \(r > 0\).

  Oletetaan, että \(f\) on tasaisesti jatkuva välillä \((0, \infty)\).
  Tällöin on olemassa \(\delta > 0\) siten, että
  \[
    \abs{f(x) - f(y)} < \delta
  \]
  kaikilla \(x,y \in (0, \infty)\) joilla \(\abs{x - y} < \delta\).
  Valitaan \(n > 1/\delta\) ja asetetaan
  \[
    x = \frac{1}{n} \quad \text{ja} \quad y = \frac{1}{n + 1}.
  \]
  Tällöin
  \[
    \abs{x - y} = \frac{1}{n(n + 1)} < \frac{1}{n} < \delta,
  \]
  mutta
  \[
    \abs{f(x) - f(y)} = \abs{n - (n + 1)} = 1 \geq \frac{1}{2}.
  \]
  Olemme löytäneet halutin ristiriidan, joten on todettava, että
  \(f\) ei ole tasaisesti jatkuva välillä \((0, \infty)\).
\end{solution}

\begin{exercise}
  Olkoon \(f \colon \RR \to \RR\) jatkuva funktio, jolle \(0 < f(x) < 1\)
  kaikilla \(x \in \RR\). Määritellään funktio \(g \colon \RR \to \RR\)
  asettamalla
  \[
    g(x) = \frac{f(x)}{1 + x^2}.
  \]
  Osoita, että funktio \(g\) saavuttaa joukossa \(\RR\) suurimman arvonsa.
\end{exercise}

\begin{solution}
  Funktio \(x \mapsto 1+x^2\) on polynomifunktiona jatkuva ja
  kaikkialla positiivinen. Koska myös \(f\) on jatkuva, osamäärän
  jatkuvuuden nojalla \(g\) on jatkuva joukossa \(\RR\).

  Olkoon \(\varepsilon > 0\) ja valitaan
  \(H = \maxset{1,1/\varepsilon}\). Kun \(\abs{x} > H\), saadaan
  \[
    \abs{g(x)} = \frac{f(x)}{1+x^2}
    < \frac{1}{1+x^2} < \frac{1}{\abs{x}} < \varepsilon.
  \]
  Täten \(\lim_{x \to \infty} g(x) = \lim_{x \to -\infty} g(x) = 0\).

  Huomataan, että \(g(0) = f(0) > 0\). Edellisten raja-arvojen nojalla
  voidaan valita \(m < 0\) ja \(M > 0\) siten, että \(\abs{g(x)} < g(0)\),
  kun \(x < m\) tai \(x > M\). Tällöin erityisesti \(g(x) < g(0)\).

  Koska \(g\) on jatkuva suljetulla ja rajoitetulla välillä \([m,M]\),
  Weierstrassin min-max-lauseen nojalla on olemassa \(s \in [m,M]\)
  siten, että
  \[
    g(s) \geq g(x) \qquad \text{kaikilla } x \in [m, M].
  \]
  Lisäksi \(0 \in [m,M]\), joten välin ulkopuolella pätee
  \[
    g(s) \geq g(0) > g(x).
  \]
  Siis \(g\) saavuttaa suurimman arvonsa pisteessä \(s\) koko joukossa
  \(\RR\).
\end{solution}

\end{document}
